\chapter{Discussion}
\label{chap:discussion}

\section{Interpretation of the primary findings}

This thesis examined whether local OCT information, once spatially aligned with fundus-controlled microperimetry, could estimate pointwise retinal sensitivity for participants with Usher syndrome who were not used to train the model. The primary results provide evidence that it can contribute useful information in this cohort. The location-plus-RETFound ridge achieved an out-of-fold MAE of 5.485~dB, compared with 7.650~dB for the location-only model and 9.346~dB for the training-mean reference. OCT-derived information was therefore associated with a 2.165~dB reduction in average absolute error beyond what could be obtained from retinal position alone. This is important because retinal position is itself informative in a degenerative retinal disease: a model that only learned that central and peripheral test locations tend to have different sensitivities could appear useful without learning any local structural relationship.

The remaining metrics support a cautious but positive interpretation. The primary combined model had an RMSE of 6.986~dB, a calibration slope of 0.978 and a mean bias of $-0.060$~dB. Its predictions were therefore not systematically shifted above or below the observed sensitivities across the cohort. At the same time, the 95\% limits of agreement extended from $-13.754$ to 13.634~dB. This range shows that individual point estimates can differ substantially from the recorded MAIA value even when the average bias is close to zero. The model should consequently be understood as an internally validated estimate of a structure--function relationship, not as an interchangeable substitute for a microperimetry measurement.

The results also demonstrate why several measures are needed. The supplementary $R^2$ of 0.564 indicates that the out-of-fold predictions captured a meaningful proportion of variation in the recorded endpoint, while MAE describes the typical size of an error on the dB scale. Bland--Altman agreement adds a different perspective by showing the range of error that could occur for an individual point. Taken together, these findings support the conclusion that the local OCT representation contains information about pointwise sensitivity in this cohort, but that this information is not sufficient to make highly precise predictions for every retinal location.

\section{Local OCT context and the effect of within-footprint pooling}

The spatial design was central to this result. A MAIA sensitivity value is produced by a stimulus with an approximately 100~$\mu$m diameter, whereas an OCT volume contains a much denser and wider set of structural measurements. Registration alone does not solve this difference in scale. The image input also has to remain connected to the retinal area represented by the functional test, which is why the three nearest B-scans and the native anchor within the stimulus footprint were defined geometrically before sensitivity was introduced.

Several A-scan positions within the physical footprint were explored during development in an attempt to represent the stimulus area more completely. In the final analyses, this within-footprint pooling was not retained because combining several nearby positions created smoother representations and produced poorer held-out performance than a single centred native anchor on each B-scan. A reasonable interpretation is that neighbouring positions can contain local structural variation that remains relevant to a pointwise endpoint; averaging them can therefore reduce the contrast between a transition region and its surroundings. This is an empirical finding from the present cohort rather than a claim that image pooling is generally inappropriate.

It is also important to distinguish this choice from the use of three B-scans. The final representation still combined the embeddings from the three nearest B-scans, using their mean and population standard deviation, because these were three immediately neighbouring cross-sectional views of the same registered functional location. Pooling across nearby positions along a single B-scan and combining these neighbouring B-scan views answer different questions. The former broadened the local anchor within the stimulus footprint and was less favourable in this study, while the latter reduced dependence on any one cross-section and retained information about variation between scans. This distinction allows the model input to remain local without treating one A-scan as a complete biological description of the stimulus area.

The primary 750~$\mu$m context crop should be interpreted in the same way. It provided surrounding retinal structure to the encoder, but its centre was fixed by the single footprint-aware anchor. The larger visual context therefore did not redefine the functional target. It allowed the representation to include neighbouring morphology that may help place a local structural change within its retinal setting, while the geometric selection rule preserved the link between the target and the OCT input.

\section{Relationship to previous structure--function studies}

The findings are consistent with the broader biological pattern described in Chapter~\ref{chap:background}: preservation of outer-retinal structure is associated with better visual function in inherited retinal disease, although the strength and form of this relationship depend on the clinical endpoint and disease stage. In the baseline RUSH2A analysis, larger ellipsoid-zone area and central retinal thickness were associated with better mean microperimetric sensitivity in participants with \textit{USH2A}-related retinal degeneration \citep{Lad2022}. The later longitudinal RUSH2A study likewise found change in OCT measures and microperimetry over four years, including a decline in mean sensitivity and in sensitivity measured at functional transition points \citep{Vincent2025}. These studies make the present result biologically plausible, but they examine cohort-level or longitudinal associations rather than an image-based prediction for each point in participants held out from training.

There is also a useful distinction between the current result and the numerical values reported in other retinal prediction studies. In geographic atrophy, Pfau et al. reported an imaging-only MAE of 4.64~dB and a lower error after patient-specific calibration \citep{Pfau2020}. Kihara et al. reported an MAE of 3.36~dB for local OCT patches in Macular Telangiectasia type 2 \citep{Kihara2019}. These studies show that structural imaging can support functional estimation, but their errors are not direct targets for this thesis. They used different diseases, functional ranges, image inputs, sampling strategies and, in some cases, data from the participant being predicted. In particular, patient-specific calibration addresses a different clinical and statistical question from prediction for a wholly unseen participant.

The most relevant comparison is therefore the analytical standard rather than the headline MAE. The present work evaluates a pointwise target within a small Usher syndrome cohort, uses an explicitly registered local input and keeps all observations from a participant together during validation. This is more demanding than a random split of retinal points because fellow eyes, repeated visits and nearby points cannot appear on both sides of the split. It is also why the result should not be weakened simply because it is less numerically favourable than values reported in conditions with different data structures. Conversely, a favourable internally held-out result does not establish transportability to another device, site or Usher syndrome subgroup.

The pointwise nature of the endpoint provides further context. Microperimetry is well suited to mapping local function, but it is a psychophysical test and individual points have greater test--retest variability than an eye-level mean. The REPEAT study in retinitis pigmentosa reported pointwise coefficients of repeatability of approximately 8.6--10.2~dB, compared with lower variability for mean sensitivity \citep{Karuntu2025}. Those values cannot be used as a direct threshold for this cohort or for a model prediction, since the study population and test protocol differ. They do, however, show why a 5.485~dB MAE should not be interpreted in the same way as an error for average sensitivity, and why repeated functional measurements would be valuable in a future validation study.

\section{What the image-derived signal may represent}

The improvement beyond the location-only model indicates that the image representation was not acting only as a proxy for the fixed MAIA grid. It does not, however, identify one structure that explains sensitivity. The OCT context may encode several related features, including the continuity of reflective outer-retinal bands, local retinal thickness, the sharpness of structural transitions and the wider pattern of degeneration around the point. These are clinically plausible sources of information because photoreceptor preservation is linked to retinal function, but the frozen representation does not provide a direct biological measurement of any one feature \citep{Tao2016,Gill2022}.

This is one reason why the result is best described as image-derived structural information rather than as proof of a particular anatomical mechanism. A point with lower recorded sensitivity may be associated with a local structural change, but it may also sit within a broader area of degeneration that is represented in the surrounding OCT context. Location, disease severity, visit and unmeasured clinical characteristics can all be related. The comparison with the location-only model helps to isolate added value beyond fixed grid position, yet it does not remove every possible source of confounding. A model trained on a larger and more heterogeneous cohort would be needed to test whether the same representation remains informative after these relationships vary more widely.

The choice of a frozen RETFound encoder was helpful in this setting because it allowed a high-dimensional retinal image representation to be used without fitting a large image model on 22 participants. This constraint was not intended to imply that a pre-trained representation is automatically appropriate for Usher syndrome. RETFound was developed and evaluated across retinal imaging tasks, and the literature shows that the transfer of an OCT model to a new disease needs to be tested rather than assumed \citep{Zhou2023RETFound,Loo2022}. In the present study, the observed improvement over both non-imaging references provides that cohort-specific test. It supports the use of the representation for this particular question, while leaving disease-specific adaptation as a question for future work.

The findings also place the role of explicit engineered features in context. The primary representation is less directly interpretable than a manually measured thickness or boundary distance, but it can retain information from the local image that may not be represented by a small set of summary variables. Conversely, the secondary random forest shows that compact structural features can perform well when verified source data are available. The two approaches should therefore be viewed as complementary. Image-derived features can be useful when consistent boundary information is unavailable, while interpretable measurements can help examine more specific structural hypotheses once their anatomical source is clear.

\section{Secondary boundary-informed findings}

The secondary analysis provided a separate view of the structure--function question in the subset for which source-mapped boundary information was available. Within this cohort, the location-plus-boundary random forest achieved the lowest MAE at 5.054~dB. The boundary-aligned RETFound model was close by MAE at 5.114~dB, while it had the lower RMSE of 6.421~dB and the higher supplementary $R^2$ of 0.602. The two models therefore appear to capture related, but not identical, structural information. The random forest can represent non-linear combinations of the local boundary measurements and retinal location, whereas the image model summarises texture and morphology from the boundary-aligned OCT context without assigning names to the source boundaries.

The fusion model did not improve on either of these models, with an MAE of 5.577~dB. This result is informative because adding more representation types does not necessarily add independent signal, particularly in a small participant cohort. Boundary features, boundary-aligned image features and retinal location can be correlated descriptions of the same disease pattern. A larger fused predictor block can then add variability without adding enough new information to improve predictions for unseen participants. The result supports the use of a model ladder in which each additional input is assessed against simpler alternatives, rather than assuming that a combined model must be stronger.

The secondary results must remain separate from the primary model comparison. Although both analyses included 22 participants, the boundary-available subset contained 50 eye-visits and 1,850 points, rather than the 78 eye-visits and 2,886 points in the full cohort. The smaller error in the secondary analysis may consequently reflect differences in eligibility, visit composition or sensitivity distribution as well as the representations and models under consideration. It cannot be interpreted as evidence that boundary features are superior to the full-cohort representation.

The neutral B0--B4 labels also constrain biological interpretation in an appropriate way. The results show that the verified source-derived intervals were useful predictors in this subset; they do not establish that a particular named retinal layer causes a change in sensitivity. Independent confirmation of the anatomical definitions and segmentation provenance would be needed before relating an individual interval to a specific outer-retinal structure. This is particularly relevant in Usher syndrome, where the relation between structural preservation, degeneration stage and measured function may vary across locations and participants \citep{Gill2022,Loo2022}.

\section{Error patterns, strengths and limitations}

The descriptive strata show where the primary model was less and more reliable within the available data. MAE was 8.787~dB for observations recorded at the $-1$~dB MAIA floor, compared with 4.737~dB above the floor. At the floor, several levels of severe functional loss can receive the same recorded value. This reduces the information available to a regression model and makes the observed target less able to distinguish between nearby levels of impairment. It also means that an estimate near the floor cannot be read as a precise measure of residual function. The lower error at year one than at baseline is descriptive only, since the visit groups differed in their numbers of observations and participating individuals.

Several features of the study strengthen the interpretation. OCT and MAIA data were paired at the same visit, and the registration was visually reviewed before pointwise modelling. The primary analysis retained all eligible matched observations rather than restricting the main question to the boundary-available cases. It also used transparent non-imaging references, a constrained representation model and nested participant-grouped validation. These decisions make the reported improvement over location less likely to arise solely from spatial co-location or from correlated observations being shared between the training and test sets.

However, internal validation does not create an external cohort. The number of participants was 22, and the apparently larger number of point rows includes correlations within eyes, visits and participants. Participant-level folds and participant-bootstrap intervals address this clustering during evaluation, but they cannot represent disease diversity that is absent from the data. The models also rely on the quality of the MAIA--OCT/SLO registration, OCT acquisition and MAIA response. A small local misregistration, poor fixation or an unobserved scan artefact can affect an individual prediction even when the overall transform has passed review.

The endpoint itself represents recorded sensitivity, not an error-free latent measure of retinal function. Fatigue, attention, learning effects and disease severity can affect microperimetry, while the device floor creates a censored part of the outcome range \citep{Pfau2021Microperimetry,Karuntu2025}. The present analyses therefore cannot separate all functional measurement variability from structural prediction error. They are also cross-sectional in their main inference: paired baseline and year-one visits contribute observations, but the study was not designed to estimate whether an OCT model detects within-person progression or a response to treatment.

The study also makes no claim about diagnosis. Its endpoint was a continuous sensitivity value at a known retinal location in people already represented within an Usher syndrome cohort. It cannot determine whether an individual has Usher syndrome, distinguish between genetic subtypes, or infer a participant's broader visual experience from a local dB estimate. These boundaries matter because a useful research model can still have a narrow intended use. Keeping that use specific makes it possible to evaluate the model against the appropriate reference---direct functional testing at the corresponding retinal point---instead of attaching wider clinical claims to its performance.

\section{Implications for research and clinical monitoring}

The appropriate implication is not that OCT can replace microperimetry. The wide individual-point agreement limits, the device floor and the absence of external validation rule out that conclusion. Instead, the study provides a reproducible, spatially explicit basis for asking when OCT carries useful information about local function. This could be valuable in research settings where microperimetry is burdensome, where an objective structural complement is needed, or where a future trial wishes to understand whether structural and functional changes move together.

For Usher syndrome studies, this question is increasingly relevant because clinical trials require endpoints that are reliable, sensitive to change and suited to the stage of disease being studied. Recent RUSH2A endpoint work emphasises the difficulty created by slow progression and variation between participants, and recommends that endpoint selection be linked to the capacity to observe change in untreated eyes \citep{Maguire2024}. The present model is not an endpoint proposal, but it suggests that registered local OCT could be investigated alongside direct functional measures rather than treated as a proxy without validation. Any future use should specify the intended decision, the acceptable uncertainty for that decision and the way in which an uncertain or failed prediction would be handled.

There may also be a practical role in supporting the interpretation of incomplete functional data. For example, an image-derived estimate could help describe the structural context of a missing or unreliable MAIA point during a research visit, provided that it is clearly labelled as an estimate rather than treated as an observed response. This would still require external evidence that the model is reliable in the relevant setting. Its value would depend on whether it improves the quality of a research assessment, not simply on whether it produces a plausible number.

Overall, the results support a positive but limited conclusion. Local OCT representations improved held-out prediction of pointwise MAIA sensitivity beyond retinal location in this Usher syndrome cohort. The smaller boundary-informed analysis further showed that source-mapped structural intervals and boundary-aligned image representations each contained useful signal, although their relative performance cannot be compared directly with the primary analysis. The work therefore supports further investigation of OCT-based functional estimation while retaining direct microperimetry as the reference functional measurement.
